%% ===================================================================================
%% Capítulo 3 -- Processo de Desenvolvimento e Ambiente
%% Fontes: CLAUDE.md, ARQUITETO.md, DESENVOLVEDORES.md. Dono: Arquiteto/Redator.
%% Este capítulo descreve o processo; as regras operacionais vivem nos arquivos de
%% governança na raiz do repositório e devem permanecer coerentes.
%% ===================================================================================
\chapter{Processo de Desenvolvimento e Ambiente}\label{cap:metodologia}

Este capítulo descreve a abordagem e as etapas do desenvolvimento (Seção~\ref{sec:abordagem}), o planejamento por fases (Seção~\ref{sec:planejamento}) e, principalmente, o processo de desenvolvimento orientado por documentação e executado por agentes especializados (Seção~\ref{sec:processo-ddd}). Encerra com o ambiente técnico (Seção~\ref{sec:ambiente}), a estratégia de verificação (Seção~\ref{sec:estrategia-vv}) e os riscos do projeto (Seção~\ref{sec:riscos}).

\section{Abordagem e Etapas do Desenvolvimento}\label{sec:abordagem}

O Vitrine é construído de forma iterativa e incremental \cite{sommerville2011}: cada incremento entrega uma fatia funcional de ponta a ponta --- da persistência no SQLite à renderização no \en{front-end} ---, verificada antes do incremento seguinte. A partir da consolidação do monólito, a especificação passa a ser escrita antes da implementação e atualizada após cada integração (Seção~\ref{sec:processo-ddd}).

O processo percorre quatro macroetapas, fundamentadas nos processos clássicos de engenharia de \en{software} \cite{sommerville2011,pressman2021}:

\begin{enumerate}
  \item \textbf{Levantamento de requisitos:} identificação das necessidades dos usuários --- professores e desenvolvedores --- e definição dos requisitos funcionais (RF) e não funcionais (RNF);
  \item \textbf{Modelagem de sistemas:} uso de UML e DBML para visualizar a lógica do sistema, incluindo diagramas de casos de uso e modelagem de dados (ER);
  \item \textbf{Prototipagem:} criação de protótipos de baixa fidelidade para validar a usabilidade do editor de Markdown e do painel administrativo antes da implementação final;
  \item \textbf{Desenvolvimento incremental (\en{sprints}):} implementação do código-fonte com Next.js e Prisma ORM, da persistência no SQLite até a renderização no \en{front-end} e a configuração de cache (ISR).
\end{enumerate}

\section{Planejamento por Fases}\label{sec:planejamento}

A consolidação do sistema em arquitetura monolítica documentada é organizada em fases sucessivas. Cada fase termina com a integração das ramificações correspondentes na \texttt{main}, a atualização da documentação e a revisão do estado dos módulos (Seção~\ref{sec:modulos}). As fases não têm datas fixas neste documento: o avanço é registrado no mural de agentes e no Capítulo~\ref{cap:implementacao}.

\begin{table}[htb]
\caption{Fases planejadas da consolidação}\label{tab:fases}
\small
\begin{tabularx}{\linewidth}{@{}L{1.1cm}L{4.2cm}YL{2.8cm}@{}}
\toprule
\textbf{Fase} & \textbf{Nome} & \textbf{Entregas principais} & \textbf{Agentes}\\
\midrule
F1 & Especificação & Documentação \LaTeX{} inicial, diagramas como código, governança e mural & \agente{Arquiteto}\\
\addlinespace
F2 & Fundação & Esquema, migração e \en{seed}; acesso a dados; autenticação; proteção de rotas & \agente{Dados}, \agente{Core}\\
\addlinespace
F3 & Conteúdo e leitura & API de conteúdo; páginas públicas com ISR; \en{sidebar} e barra superior; renderização Markdown com \LaTeX{} & \agente{Core}, \agente{UI}\\
\addlinespace
F4 & Edição e painel & Editor com pré-visualização; gestão da árvore (criar, mover, ordenar, excluir) & \agente{UI}, \agente{Core}\\
\addlinespace
F5 & Personalização & Identidade e paleta por variáveis CSS; revalidação de \en{layout} & \agente{UI}, \agente{Core}, \agente{Dados}\\
\addlinespace
F6 & Verificação e validação & Testes automatizados, desempenho, segurança, acessibilidade e usabilidade & Todos\\
\addlinespace
F7 & Consolidação da documentação & Relato de implementação por módulo e situação da verificação & \agente{Arquiteto}\\
\bottomrule
\end{tabularx}
\fonte{Elaborado pelos autores.}
\end{table}

\section{Processo de Desenvolvimento: Documentation-Driven Development com Agentes}\label{sec:processo-ddd}

\subsection{Princípios}\label{sec:principios-ddd}

O processo segue sete princípios, que também constam dos arquivos de governança do repositório.

\begin{enumerate}[label=\textbf{PR\arabic*.}, leftmargin=1.6cm]
  \item \textbf{A documentação é a fonte da verdade.} Em caso de divergência entre código e documentação, a divergência é tratada como defeito de um dos dois e resolvida explicitamente; nunca é ignorada.
  \item \textbf{A especificação precede a implementação.} Nenhuma funcionalidade nova é implementada sem que exista, na documentação, o requisito, o contrato e o critério de aceite correspondentes.
  \item \textbf{Escritor único da documentação.} Somente o agente Arquiteto edita a documentação (\arq{doc/}) e os arquivos de governança; os demais agentes propõem mudanças por meio do mural.
  \item \textbf{Domínios com dono.} Cada diretório do código pertence a um agente; alterações fora do domínio são solicitadas ao dono.
  \item \textbf{Código apenas em ramificações.} É vedado alterar código ou configuração diretamente na \texttt{main}; a integração ocorre por \en{merge} autorizado pelo Arquiteto. Cada agente trabalha em um \en{worktree} próprio, e a cópia principal permanece na \texttt{main}.
  \item \textbf{O mural é o único canal de coordenação e só é editado na \texttt{main}} (na cópia principal). Isso serializa a comunicação e evita conflitos de fusão.
  \item \textbf{Rastreabilidade por identificadores.} Requisitos, casos de uso, regras, decisões e módulos possuem identificadores estáveis, citados nos \en{commits}, nos pedidos de integração e nos testes.
\end{enumerate}

\subsection{Papéis e Domínios}\label{sec:papeis}

O time é composto por quatro agentes. A Tabela~\ref{tab:papeis} resume a responsabilidade e o domínio de cada um; a distribuição de diretórios é normativa e está detalhada em \arq{DESENVOLVEDORES.md}.

\begin{table}[htb]
\caption{Papéis, domínios e diretórios}\label{tab:papeis}
\small
\begin{tabularx}{\linewidth}{@{}L{2.6cm}YL{5.3cm}@{}}
\toprule
\textbf{Agente} & \textbf{Responsabilidade} & \textbf{Diretórios e arquivos}\\
\midrule
\agente{Arquiteto} (Redator Técnico e \en{Merge Master}) & Governa a fonte da verdade; revisa pedidos de integração; autoriza e executa \en{merges}; atualiza a documentação; mantém o mural limpo & \arq{doc/}, \arq{CLAUDE.md}, \arq{ARQUITETO.md}, \arq{DESENVOLVEDORES.md}, \arq{MURAL_DE_AGENTES.md} (limpeza)\\
\addlinespace
\agente{Dados} & Modela o esquema e as rotinas de acesso a dados; migrações e \en{seed} & \arq{prisma/}, \arq{src/lib/db.js}, \arq{src/lib/content-tree.js}\\
\addlinespace
\agente{Core} & Lógica de servidor: API, autenticação, proteção de rotas, revalidação de cache; roteamento do App Router & \arq{src/app/api/}, \arq{src/proxy.js}, \arq{src/lib/auth*.js}, \arq{src/lib/validators.js}, \arq{src/lib/revalidate.js}, arquivos de rota (\arq{page.js}, \arq{layout.js}, \arq{route.js})\\
\addlinespace
\agente{UI} & Interface: componentes, estilo global, editor e pipeline Markdown & \arq{src/components/}, \arq{src/app/globals.css}, \arq{src/lib/theme.js}\\
\bottomrule
\end{tabularx}
\fonte{Elaborado pelos autores, com base em \arq{DESENVOLVEDORES.md} e \arq{ARQUITETO.md}.}
\end{table}

\subsection{Fontes da Verdade}\label{sec:fontes-verdade}

A Tabela~\ref{tab:fontes-verdade} relaciona cada tipo de informação ao artefato que o define. Quem implementa deve consultar esses arquivos antes de escrever código e citar os identificadores (RF, RNF, RN, UC, ADR, D) no pedido de integração.

\begin{table}[htb]
\caption{Fontes da verdade do projeto}\label{tab:fontes-verdade}
\small
\begin{tabularx}{\linewidth}{@{}L{3.3cm}YL{3.5cm}@{}}
\toprule
\textbf{Informação} & \textbf{Artefato (caminho)} & \textbf{Responsável}\\
\midrule
Requisitos, regras e casos de uso & \arq{doc/tex/capitulos/cap4_requisitos.tex}; \arq{doc/diagrams/src/uc-visao-geral.puml} & \agente{Arquiteto}\\
\addlinespace
Arquitetura, API, rotas e cache & \arq{doc/tex/capitulos/cap5_arquitetura.tex}; diagramas \arq{arq-*} e \arq{seq-*} & \agente{Arquiteto}\\
\addlinespace
Modelo de dados & Implementação: \arq{prisma/schema.prisma}; documentação: \arq{doc/diagrams/src/er-vitrine.dbml} e \arq{doc/tex/apendices/apendice_a_dicionario_dados.tex} & \agente{Dados} (código); \agente{Arquiteto} (documentação)\\
\addlinespace
Decisões de projeto & \arq{doc/tex/apendices/apendice_e_adr.tex} & \agente{Arquiteto}\\
\addlinespace
Estado de cada módulo & \arq{doc/tex/capitulos/cap6_implementacao.tex} e mapa de pendências (Apêndice~\ref{ap:pendencias}) & \agente{Arquiteto}\\
\addlinespace
Processo e regras de trabalho & \arq{CLAUDE.md}, \arq{ARQUITETO.md}, \arq{DESENVOLVEDORES.md} & \agente{Arquiteto}\\
\addlinespace
Comunicação e pedidos de integração & \arq{MURAL_DE_AGENTES.md} (somente na \texttt{main}) & Todos\\
\bottomrule
\end{tabularx}
\fonte{Elaborado pelos autores.}
\end{table}

\subsection{Fluxo de Git e Uso do Mural}\label{sec:fluxo-git}

O ciclo completo está na Figura~\ref{fig:fluxo-ddd-agentes}. Em resumo: o Arquiteto especifica e publica avisos no mural; cada desenvolvedor cria, no primeiro passo, um \en{worktree} próprio com uma ramificação a partir da \texttt{main} e implementa nele, conforme a especificação; ao terminar, vai à cópia principal, que permanece na \texttt{main}, apenas para registrar o pedido de integração no mural; o Arquiteto compara a ramificação com a \texttt{main}, autoriza ou veta, e após o \en{merge} atualiza a documentação e limpa o mural.

\vitrinefigura[tipo=diagrama, arquivo=fluxo-ddd-agentes, agente=Arquiteto, largura=0.95,
  origem={doc/diagrams/src/fluxo-ddd-agentes.puml},
  descricao={Diagrama de atividades com raias (Arquiteto, Desenvolvedor) mostrando especificação, implementação, pedido no mural, revisão, merge ou veto e atualização da documentação.}]
  {Ciclo de \en{Documentation-Driven Development} com agentes}{fig:fluxo-ddd-agentes}

O mecanismo de fusão assíncrona por ramificações e a regra de que o mural só é editado na \texttt{main} formam um protocolo simples: como cada pedido é um \en{commit} na \texttt{main}, o histórico do Git registra quem pediu o quê e quando (Figura~\ref{fig:fluxo-git-mural}).

\textbf{Cópia de trabalho própria.} Os agentes executam ao mesmo tempo e não se veem. Se dividissem a mesma pasta, a troca de ramificação de um deslocaria a ramificação ativa do outro, e um \en{commit} cairia na ramificação errada --- o que ocorreu em 02/10/2026, quando um \en{commit} de código foi parar na \texttt{main} e outro, na ramificação de um colega (risco R15). Por isso cada agente cria, como primeiro passo, um \en{worktree} (\texttt{git worktree add ../vitrine-<agente> -b <ramificação> main}); a pasta principal fica sempre na \texttt{main}, serve só ao mural e ao Arquiteto, e ninguém troca a ramificação dela; confere-se \texttt{git branch -{}-show-current} antes de cada \en{commit}; e o Arquiteto remove o \en{worktree} depois do \en{merge}. Um \en{commit} de código que chegue à \texttt{main} por engano não é apagado do histórico: é revisado, mantido e registrado como exceção no \en{merge} correspondente.

\vitrinefigura[tipo=diagrama, arquivo=fluxo-git-mural, agente=Arquiteto, largura=0.95,
  origem={doc/diagrams/src/fluxo-git-mural.puml},
  descricao={Diagrama de sequência entre Desenvolvedor, ramificação de feature, main, MURAL\_DE\_AGENTES.md e Arquiteto, incluindo os caminhos de aprovação e de veto.}]
  {Fluxo Git e solicitação de integração pelo mural}{fig:fluxo-git-mural}

\subsection{Critérios de Aceite de uma Integração}\label{sec:criterios-aceite}

Ao revisar um pedido, o Arquiteto verifica, no mínimo:

\begin{enumerate}
  \item a ramificação parte da \texttt{main} atualizada e altera apenas arquivos do domínio do autor (ou declara as exceções no pedido);
  \item o comportamento implementado corresponde aos requisitos, às regras e aos contratos citados (identificadores RF, RNF, RN, UC);
  \item alterações de esquema, de API e de dependências estão declaradas no pedido e são compatíveis com a documentação --- ou acompanham um pedido de atualização da documentação;
  \item \texttt{npm run lint} e \texttt{npm run build} passam, e os testes aplicáveis foram executados;
  \item nenhum segredo, arquivo gerado ou banco de dados é versionado;
  \item há evidências visuais (capturas de tela nos caminhos previstos) quando a mudança é visível ao usuário;
  \item a ramificação não altera o mural nem a documentação (\arq{doc/});
  \item os critérios de segurança aplicáveis foram atendidos: validação de entrada no servidor, sessão verificada no servidor e ausência de HTML bruto no conteúdo.
\end{enumerate}

Se algum critério falha, o Arquiteto não realiza o \en{merge} e responde no mural, apontando o que deve ser corrigido.

\subsection{Rastreabilidade e Identificadores}\label{sec:rastreabilidade}

Todos os artefatos usam identificadores estáveis: \texttt{UC} (casos de uso), \texttt{RF} e \texttt{RNF} (requisitos), \texttt{RN} (regras de negócio), \texttt{ADR} (decisões aceitas) e \texttt{D} (identificadores das decisões iniciais, hoje todas registradas nas ADR), \texttt{M} (módulos) e \texttt{CT} (casos de teste). A Seção~\ref{sec:matriz-rastreabilidade} relaciona requisitos, casos de uso e módulos, e o Capítulo~\ref{cap:testes} relaciona requisitos e casos de teste. As referências cruzadas no documento são hiperligações: clicar em \req{RF01} leva à definição do requisito.

\subsection{Convenções de Ramificações, \emph{Commits} e Idioma}

As ramificações seguem o padrão \arq{tipo/dominio-tema} --- por exemplo, \arq{feat/dados-schema-inicial} ou \arq{fix/core-sessao-expirada} --- e as mensagens de \en{commit}, o padrão \en{Conventional Commits} \cite{conventionalcommits} em português, com os identificadores afetados em rodapé (por exemplo, \texttt{Refs: RF05, RN03}). Identificadores de código e nomes de arquivos são em inglês; comentários, textos de interface e documentação são em português do Brasil.

\subsection{Uso de Agentes de Inteligência Artificial e Responsabilidade}\label{sec:ia-autoria}

Os agentes de inteligência artificial são ferramentas de apoio sob supervisão humana: não assumem responsabilidade pelo produto. O Engenheiro Chefe (autor do projeto) é o responsável pela revisão do código e da documentação que integram o produto, pelas decisões de projeto e pela veracidade das referências. O registro das interações relevantes --- pedidos, vetos e decisões --- é preservado no histórico do Git, o que permite auditar como cada parte do sistema foi produzida.

\section{Ambiente e Ferramentas}\label{sec:ambiente}

A Tabela~\ref{tab:stack} lista as tecnologias e as versões efetivamente declaradas no repositório, conferidas em \arq{package.json} e em \arq{node_modules}; os valores foram atualizados após o merge M-001 (30/09/2026) e revistos até o merge M-011 (03/10/2026).

\begin{table}[htb]
\caption{Tecnologias e versões do projeto}\label{tab:stack}
\small
\begin{tabularx}{\linewidth}{@{}L{4.0cm}L{3.0cm}Y@{}}
\toprule
\textbf{Tecnologia} & \textbf{Versão} & \textbf{Papel}\\
\midrule
Node.js & $\geq$ 20.9 (24 no desenvolvimento e na imagem) & \en{Runtime} exigido pelo Next.js 16. O \arq{jsdom} dos testes de interação (merge M-011) exige 22.22.2, 24.15 ou superior: com o 22.22.0 o npm só avisa (\texttt{EBADENGINE}) e os testes passam\\
Next.js & 16.3.7 & \en{Framework} web (App Router, ISR, Route Handlers, Proxy)\\
React e React DOM & 19.2.8 & Biblioteca de interface\\
React Compiler (plugin Babel) & 1.0.0 & Memoização automática (\texttt{reactCompiler: true})\\
Tailwind CSS & 4.3.3 (instalada; \texttt{\^{}4}) & Estilização por utilitários\\
Prisma ORM e \arq{@prisma/client} & 7.10.0 & Acesso a dados e migrações; configuração em \arq{prisma.config.ts}; cliente gerado em \arq{src/generated/prisma} (gerador \texttt{prisma-client})\\
\arq{@prisma/adapter-better-sqlite3} e \arq{better-sqlite3} & 7.10.0 e 12.11.1 & \en{Driver adapter} do Prisma para SQLite (módulo nativo) e o próprio SQLite em arquivo (\adr{ADR-014})\\
\arq{prisma} (CLI) e \arq{dotenv} & 7.10.0 e 18.0.5 (\texttt{dependencies}) & Migrações e carregamento do \arq{.env} pela CLI; necessários na imagem de produção (\texttt{prisma migrate deploy})\\
Vitest & 5.0.3 (desenvolvimento) & Testes de unidade e de integração com banco SQLite isolado (\texttt{npm test}; Seção~\ref{sec:testes-banco})\\
\arq{jsdom}, \arq{@testing-library/react} e \arq{@testing-library/dom} & 30.1.1, 16.3.3 e 10.4.2 (desenvolvimento; MIT) & Testes de interação dos Componentes Cliente (cliques, estado, \texttt{localStorage}): cada arquivo de teste escolhe o ambiente com \texttt{// @vitest-environment jsdom}, sem mudar a configuração do Vitest; 0~byte no \en{bundle} (merge M-011)\\
\arq{@playwright/test} & 1.63.0 (desenvolvimento; Apache-2.0) & Testes de ponta a ponta no navegador, sobre o \en{build} de produção (\texttt{npm run test:e2e}; \decisao{D-12}; merge M-010). Os navegadores ficam fora do repositório e da imagem; os três pacotes (\texttt{@playwright/test}, \texttt{playwright} e \texttt{playwright-core}) entram na imagem como qualquer dependência de desenvolvimento, cerca de 18~MiB em \arq{node_modules}\\
\arq{tsx} & 4.23.15 (desenvolvimento) & Execução do \en{seed} (\arq{prisma/seed.js}) com o cliente gerado em TypeScript\\
better-auth & 1.7.6 & Autenticação e sessões (e-mail e senha, cadastro público desabilitado, limitação de tentativas)\\
react-markdown & 10.1.0 & Renderização de Markdown em React\\
remark-math e rehype-katex & 6.0.0 e 7.0.1 & Delimitadores e renderização de fórmulas\\
remark-gfm & 4.0.1 & Tabelas, tachado, listas de tarefas, links automáticos e notas de rodapé (RF11; \decisao{D-08}; merge M-009)\\
rehype-highlight, \arq{lowlight} e \arq{highlight.js} & 7.0.2, 3.3.0 e 11.11.2 & Realce de sintaxe por classes, no servidor e no cliente (RF10; \decisao{D-07}; merge M-009; o \arq{highlight.js} é BSD-3-Clause, os demais são MIT). No merge M-011, o \arq{highlight.js} (fixado em \texttt{\~{}11.11.2}) e o \arq{lowlight} passaram a dependências diretas, com LaTeX, MATLAB, Julia, Haskell e Fortran (42 linguagens; mais 6,3~kB comprimidos)\\
react-hook-form & 7.89.0 (MIT) & Formulário do editor, com validação por campo e foco no primeiro erro, sobre os esquemas do Zod por um resolvedor próprio (M09; merge M-011). Só entra no código do editor, carregado sob demanda (12,4~kB comprimidos)\\
KaTeX & 0.16.47 & Composição de fórmulas; declarado diretamente no merge M-009, pois o \arq{MarkdownRenderer} importa a folha de estilos do pacote\\
Zod & 4.6.5 (\texttt{dependencies}) & Validação de entrada da API e do comando do administrador (\decisao{D-11}); declarado diretamente no merge M-004 (já era dependência do better-auth)\\
ESLint e \arq{eslint-config-next} & 9.x e 16.3.7 & Análise estática\\
\LaTeX{} (MiKTeX) e \texttt{latexmk} & pdfTeX 4.19 & Composição desta documentação\\
PlantUML e Graphviz & a instalar & Renderização dos diagramas (\arq{doc/diagrams/})\\
Git & --- & Controle de versão e fluxo por ramificações\\
Docker Engine e Docker Compose & versão do servidor & Implantação por imagem em VPS tradicional (\adr{ADR-011})\\
\bottomrule
\end{tabularx}
\fonte{Elaborado pelos autores, a partir de \arq{package.json} e das instalações locais.}
\end{table}

\section{Estratégia de Verificação e Validação}\label{sec:estrategia-vv}

A verificação (o sistema foi construído corretamente?) e a validação (foi construído o sistema correto?) serão realizadas em quatro frentes: (i) testes automatizados de unidade, integração e ponta a ponta sobre os casos de uso; (ii) medições de desempenho e de comportamento do cache ISR; (iii) verificações de segurança e acessibilidade orientadas pelo OWASP e pelo WCAG; e (iv) avaliação de usabilidade com usuários, com o questionário SUS. O plano detalhado, com critérios e casos de teste, está no Capítulo~\ref{cap:testes}.

\section{Riscos do Projeto}\label{sec:riscos}

A Tabela~\ref{tab:riscos} registra os principais riscos identificados, com probabilidade (P) e impacto (I) qualitativos --- baixo (B), médio (M) ou alto (A) --- e as mitigações adotadas.

\begin{longtable}{@{}L{1.0cm}L{5.0cm}L{0.7cm}L{0.7cm}L{6.9cm}@{}}
\caption{Registro de riscos}\label{tab:riscos}\\
\toprule
\textbf{ID} & \textbf{Risco} & \textbf{P} & \textbf{I} & \textbf{Mitigação}\\
\midrule
\endfirsthead
\toprule
\textbf{ID} & \textbf{Risco} & \textbf{P} & \textbf{I} & \textbf{Mitigação}\\
\midrule
\endhead
\bottomrule
\endfoot
R01 & O sistema de arquivos efêmero de plataformas \en{serverless} impede a persistência do arquivo SQLite & B & A & Mitigado pela decisão \decisao{D-04}: implantação por imagem Docker em VPS com volume persistente (\adr{ADR-011}); plataformas \en{serverless} não são alvo desta versão\\
R02 & Divergência entre código e documentação (deriva) & M & A & Princípios PR1 e PR2; critérios de aceite; Arquiteto como revisor único\\
R03 & Conflitos de fusão em \arq{package-lock.json} quando várias ramificações adicionam dependências & A & B & Novas dependências declaradas no pedido; o Arquiteto regenera o \en{lockfile} ao integrar\\
R04 & Mudanças incompatíveis nas dependências (Next.js 16, Prisma 7, better-auth 1.x) em relação ao conhecimento dos agentes & M & M & Versões fixadas no \en{lockfile}; consulta obrigatória à documentação local e às \en{skills} do Prisma antes de codificar\\
R05 & Injeção de HTML/\en{script} (XSS) por meio do Markdown & M & A & Não habilitar HTML bruto; testes de segurança; revisão de dependências de Markdown\\
R06 & Cores escolhidas com contraste insuficiente & M & M & Aviso de contraste no painel (WCAG 2.1 AA); valores padrão acessíveis\\
R07 & Limite de concorrência de escrita do SQLite & B & M & Modo WAL; poucas escritas; transações curtas\\
R08 & Agentes introduzem código fora do escopo ou sem testes & M & M & Domínios com dono; critérios de aceite; veto fundamentado no mural\\
R09 & Perda de dados por exclusão em cascata acidental & M & A & Confirmação explícita; rotina de \en{backup} documentada (Apêndice~\ref{ap:implantacao})\\
R10 & Escopo maior do que a capacidade de entrega do time & M & A & Priorização Essencial, Importante e Desejável; fases com entregas incrementais\\
R11 & Vazamento de rascunhos (\texttt{DRAFT}) por consulta pública sem filtro de estado, na página, na \en{sidebar}, na geração estática ou no mapa do site & M & A & Regra RN15; filtro centralizado em \arq{db.js}; caso de teste CT31; revisão de todo \en{endpoint} ou página nova\\
R12 & Perda do volume do SQLite na VPS (falha de disco ou do provedor) & M & A & Rotina de cópia de segurança com \texttt{VACUUM INTO} e cópias fora da VPS (Apêndice~\ref{ap:implantacao}); teste de restauração (CT30)\\
R13 & Alertas de vulnerabilidade do \texttt{npm audit} em dependências transitivas do Prisma 7 (\arq{mysql2} e \arq{deepmerge-ts}), cuja correção automática rebaixaria o Prisma para a versão 6 & B & M & Alertas avaliados e aceitos (Seção~\ref{sec:arq-npm-audit}, \adr{ADR-019}): os pacotes não são exercidos pelo Vitrine; \texttt{npm audit} a cada entrega e nova avaliação quando o Prisma publicar versão corrigida\\
R14 & Perda do acesso ao painel por senha esquecida do administrador único, sem envio de e-mail nem redefinição & B & A & Guardar a senha com segurança; evolução prevista de um comando de redefinição (Seção~\ref{sec:trabalhos-futuros}); \en{backup} do banco (Apêndice~\ref{ap:implantacao})\\
R15 & Agentes simultâneos na mesma cópia de trabalho: a troca de ramificação de um desloca a do outro, e \en{commits} caem na ramificação errada (ocorrido em 02/10/2026, com um \en{commit} de código na \texttt{main}) & A & M & \en{Worktree} próprio por agente desde o primeiro passo (PR5; Seção~\ref{sec:fluxo-git}); cópia principal sempre na \texttt{main}; conferir a ramificação antes de cada \en{commit}; o Arquiteto remove o \en{worktree} após o \en{merge}\\
\end{longtable}
\fonte{Elaborado pelos autores.}
